\chapter{Data Warehousing, Modelado Dimensional y Arquitectura OLAP}
\label{cap:data_warehouse_modelado_dimensional}

\section{Marco Conceptual y Te\'orico (Curr\'iculo Universitario)}

El Data Warehouse corporativo almacena datos hist\'oricos integrados desde m\'ultiples fuentes operacionales heterog\'eneas para alimentar herramientas anal\'iticas de Business Intelligence (BI).

\subsection{Comparativa OLTP vs OLAP}
\begin{table}[h!]
\centering
\small
\begin{tabularx}{\textwidth}{lXX}
\toprule
\textbf{Dimensi\'on} & \textbf{Sistemas OLTP (Operacionales)} & \textbf{Sistemas OLAP (Anal\'iticos)} \\
\midrule
Prop\'osito & Procesamiento de transacciones en tiempo real & Soporte a decisiones y an\'alisis hist\'orico \\
Dise\~no & Altamente normalizado (3FN / BCNF) & Desnormalizado (Estrella / Copo de Nieve) \\
Patr\'on I/O & Millones de escrituras cortas (\texttt{INSERT/UPDATE}) & Consultas complejas de lectura masiva \\
M\'etrica & Rendimiento de transacciones por segundo (TPS) & Tiempo de respuesta de consultas anal\'iticas \\
\bottomrule
\end{tabularx}
\caption{Diferencias estructurales entre sistemas OLTP y OLAP}
\end{table}

\subsection{Metodolog\'ia Kimball: Hechos y Dimensiones}
\begin{itemize}
    \item \textbf{Tablas de Hechos (\textit{Fact Tables}):} Centralizan los eventos cuantitativos del negocio (montos, cantidades, duraciones). Contienen claves for\'aneas hacia las dimensiones y medidas aditivas, semi-aditivas o no aditivas.
    \item \textbf{Tablas de Dimensiones (\textit{Dimension Tables}):} Proporcionan el contexto descriptivo (\textit{qui\'en, qu\'e, d\'onde, cu\'ando}).
    \item \textbf{Dimensiones de Cambio Lento (SCD):}
    \begin{itemize}
        \item \textit{SCD Tipo 1:} Sobrescribe el valor anterior (sin historial).
        \item \textit{SCD Tipo 2:} Crea un nuevo registro con fechas de vigencia (\texttt{start\_date}, \texttt{end\_date}, \texttt{is\_current}).
        \item \textit{SCD Tipo 3:} Almacena el valor anterior en una columna separada.
    \end{itemize}
\end{itemize}

---

\section{Casos de Estudio: Operaciones OLAP (CUBE y ROLLUP)}

\begin{lstlisting}[style=sqlstyle, caption={Agregaci\'on multidimensional jer\'arquica con ROLLUP}]
-- Consulta anal\'itica con totales y subtotales autom\'aticos por jerarqu\'ia temporal
SELECT
    EXTRACT(YEAR FROM o.order_date)::INTEGER AS sales_year,
    EXTRACT(QUARTER FROM o.order_date)::INTEGER AS sales_quarter,
    p.category_id,
    SUM(od.unit_price * od.quantity * (1 - od.discount))::NUMERIC(14, 2) AS total_revenue
FROM orders AS o
INNER JOIN order_details AS od ON o.order_id = od.order_id
INNER JOIN products AS p ON od.product_id = p.product_id
GROUP BY
    ROLLUP (
        EXTRACT(YEAR FROM o.order_date),
        EXTRACT(QUARTER FROM o.order_date),
        p.category_id
    )
ORDER BY
    sales_year NULLS FIRST,
    sales_quarter NULLS FIRST,
    p.category_id NULLS FIRST;
\end{lstlisting}

---

\section{Taller de Consolidaci\'on del Cap\'itulo}

\begin{businessproblem}
Dise\~nar e implementar el esquema f\'isico de Data Warehouse para el \'area de Recursos Humanos de una multinacional, modelando la dimensi\'on de tiempo, departamento y empleados con claves subrogadas, y la tabla de hechos con costos totales y conteo de personal.
\end{businessproblem}

\subsection{Soluci\'on T\'ecnica DDL Dimensional}

\begin{lstlisting}[style=sqlstyle, caption={Esquema estrella DDL para Data Warehouse de RRHH}]
-- 1. Dimensi\'on Empleado (Con Clave Subrogada)
CREATE TABLE dim_employee (
    employee_key SERIAL PRIMARY KEY,
    employee_id INTEGER NOT NULL,
    full_name VARCHAR(100) NOT NULL,
    job_title VARCHAR(50) NOT NULL,
    salary NUMERIC(10, 2) NOT NULL,
    valid_from DATE NOT NULL DEFAULT CURRENT_DATE,
    valid_to DATE,
    is_current BOOLEAN DEFAULT TRUE
);

-- 2. Dimensi\'on Departamento
CREATE TABLE dim_department (
    department_key SERIAL PRIMARY KEY,
    department_id INTEGER NOT NULL,
    department_name VARCHAR(60) NOT NULL,
    city VARCHAR(50) NOT NULL,
    country VARCHAR(50) NOT NULL
);

-- 3. Dimensi\'on Tiempo
CREATE TABLE dim_date (
    date_key INTEGER PRIMARY KEY, -- Formato entero YYYYMMDD
    full_date DATE NOT NULL UNIQUE,
    month_number SMALLINT NOT NULL,
    month_name VARCHAR(15) NOT NULL,
    quarter_number SMALLINT NOT NULL,
    year_number SMALLINT NOT NULL
);

-- 4. Tabla de Hechos: N\'omina y Planilla
CREATE TABLE fact_payroll (
    payroll_fact_id BIGSERIAL PRIMARY KEY,
    employee_key INTEGER NOT NULL REFERENCES dim_employee(employee_key),
    department_key INTEGER NOT NULL REFERENCES dim_department(department_key),
    date_key INTEGER NOT NULL REFERENCES dim_date(date_key),
    total_cost NUMERIC(12, 2) NOT NULL,
    headcount_contribution SMALLINT DEFAULT 1
);

-- \'indices B-Tree en claves for\'aneas para acelerar JOINs dimensionales
CREATE INDEX idx_fact_payroll_emp ON fact_payroll (employee_key);
CREATE INDEX idx_fact_payroll_dept ON fact_payroll (department_key);
CREATE INDEX idx_fact_payroll_date ON fact_payroll (date_key);
\end{lstlisting}

\begin{engtip}
\begin{itemize}
    \item \textbf{Claves Subrogadas (\textit{Surrogate Keys}):} Se emplean enteros autoincrementales (\texttt{employee\_key}) en lugar de claves de negocio operacionales (\texttt{employee\_id}) para permitir el seguimiento hist\'orico de cambios mediante SCD Tipo 2 e independizar el almac\'en anal\'itico de los sistemas transaccionales fuente.
\end{itemize}
\end{engtip}
